\documentclass[runningheads]{llncs}

\usepackage[T1]{fontenc}
\usepackage{lmodern}

\usepackage{graphicx}
\usepackage{xcolor}
\usepackage{booktabs}
\usepackage{url}
\usepackage[hidelinks]{hyperref}
\usepackage{array}



\newcolumntype{L}[1]{>{\raggedright\arraybackslash}p{#1}}

\begin{document}

\title{A Drift-Aware MLOps Framework for Multi-Model Serving: A Network Intrusion Detection Case Study}

\titlerunning{A Drift-Aware MLOps Framework for Multi-Model Serving}

\author{Hoang Tran-Nguyen-Viet\inst{1} \and
Thai Bui-Ngoc\inst{1} \and
Tuan Le-Anh\inst{1} \and
Bao Nguyen-Van\inst{1} \and
Quan Le-Trung\inst{1}
}

\authorrunning{Hoang Tran-Nguyen-Viet et al.}

\institute{Faculty of Computer Networks and Communications,\\
University of Information Technology, VNU-HCM,\\Ho Chi Minh City, Vietnam\\
\email{\{23520541,23521412\}@gm.uit.edu.vn, \{tuanla,baonv,quanlt\}@uit.edu.vn}}

\maketitle

\begin{abstract}
Input-distribution drift is useful monitoring evidence, but it does not directly establish predictive degradation. We present a drift-aware MLOps prototype for tenant-owned tabular models, in which a PostgreSQL-backed Native Registry links model versions, serving metadata, inference logs, reference data, and monitoring reports. The contribution is an explicit separation of distribution alerts, labeled performance evidence, and human-reviewed lifecycle actions, rather than a new classifier or drift detector. A controlled NIDS experiment uses a fixed XGBoost model trained on BENIGN and DDoS, disjoint data pools, 120 evaluation windows, two window sizes, and four alert thresholds. Changing attack prevalence produces strong drift while macro-F1 remains near one. Replacing one quarter of attack samples with held-out attack families reduces mean attack recall to 0.7533 for 1,000-sample windows, although a corrected KS monitor alerts in only one of five windows at a drift-share threshold of 0.30. These results characterize a monitoring limitation under controlled shifts; they do not establish natural drift, retraining effectiveness, or production-scale serving.
\keywords{MLOps \and Data Drift \and Multi-model Serving \and Network Intrusion Detection}
\end{abstract}

\section{Introduction}
Network intrusion detection encounters changing traffic and attack families, while labels needed to measure predictive quality may arrive late. Generalization beyond the training environment is a longstanding concern~\cite{sommer2010closedworld}. An input-distribution change does not by itself establish predictive degradation; an absent alert does not certify continued accuracy.

MLOps coordinates models, data, deployment, and monitoring~\cite{kreuzberger2023mlops}. Our question is how to associate evidence with a tenant and model version while keeping update decisions explicit. A Native Registry and replaceable adapters implement this separation. This is a design choice, not a new learning algorithm or a capability claimed to be absent from existing platforms.

We contribute a tenant/model lifecycle design separating evidence from reviewed actions, and a fixed-model NIDS evaluation with stable controls and threshold/window-size sensitivity. We ask: \emph{under which constructed changes does drift agree or disagree with degradation?} Multi-model serving is the architectural scope; the experiment evaluates one classifier, not multi-tenant scalability.

\section{Related Work}
Kreuzberger et al.~\cite{kreuzberger2023mlops} synthesize MLOps components and workflows. Technical-debt and engineering studies emphasize data dependencies and the joint evolution of data, code, and models~\cite{sculley2015debt,amershi2019software}. Subramanya et al.~\cite{subramanya2022devops} discuss MLOps through electricity forecasting, while Renggli et al.~\cite{renggli2021data} examine how data quality affects the lifecycle. These works motivate explicit evidence and metadata management; they do not establish superiority of our implementation.

TFX integrates data validation, training, and serving~\cite{baylor2017tfx}; MLflow supports tracking, reproducibility, and deployment interfaces~\cite{mlflow}; Clipper studies low-latency prediction serving~\cite{crankshaw2017clipper}. Recent studies examine self-adaptation for sustainable MLOps~\cite{Bhatt2025HarmonE} and empirically compare workflow frameworks~\cite{MarcosMercade2026MLOpsFrameworks}. We use existing tools as adapters and evaluate a narrower monitoring question, without a cross-platform performance comparison.

Concept-drift surveys distinguish changes in distributions and predictive relationships~\cite{gama2014survey,webb2016drift}. Continual-learning NIDS methods such as SSF~\cite{Zhang2025SSFNIDS} additionally evaluate model updates. Our model remains fixed: we measure input-distribution changes and labeled performance separately, without identifying concept drift or evaluating an adaptation algorithm.

\section{Framework and Monitoring Contract}
\label{sec:framework}
\paragraph{State ownership and execution.}
The Django control plane stores tenant ownership, model versions, artifact references, readiness metadata, and drift results in PostgreSQL. FastAPI resolves a model record and invokes a loader. Endpoint metadata alone does not prove runtime readiness. Imported artifacts and training outputs share this boundary. MLflow, object storage, and workflow engines are supporting adapters, not owners of product state.

Events carry tenant/model identifiers, features, predictions, timestamps, and optional confidence. Redpanda and a consumer separate persistence from synchronous inference; Evidently provides monitoring. The offline experiment does not establish delivery completeness, failure recovery, or correctness under concurrent updates.

\begin{figure}[t]
\centering
\small
\fbox{\parbox{0.93\linewidth}{\centering
\textbf{Control plane: tenant/model Native Registry}\\
Artifact references, deployability, reference identity, monitoring results, review state}}
\par\smallskip
$\downarrow$\quad model resolution and configuration\quad$\uparrow$\quad monitoring evidence
\par\smallskip
\fbox{\parbox{0.23\linewidth}{\centering Serving\\prediction + log event}}
$\longrightarrow$
\fbox{\parbox{0.23\linewidth}{\centering Broker/consumer\\production records}}
$\longrightarrow$
\fbox{\parbox{0.23\linewidth}{\centering Monitor\\reference vs. window}}
\caption{Logical responsibilities. A distribution alert feeds review; it does not authorize an automatic model replacement. The experiment in Section~\ref{sec:protocol} isolates the monitoring comparison from this deployed pipeline.}
\label{fig:architecture}
\end{figure}

\paragraph{Distribution evidence.}
For tenant $u$ and deployable model version $i$, let $\hat y=f_{u,i}(x)$, with reference $D_{\mathrm{ref}}^{u,i}$ and current window $D_{\mathrm{cur}}^{u,i}$. A report must retain this model/reference/window association to make the evidence interpretable. For $d$ numeric features, a two-sample test produces statistics $\delta_j$ and p-values $p_j$. In the experiment, Benjamini--Hochberg (BH) adjustment~\cite{bh1995} yields $q_j$ within each window, and
\begin{equation}
I_j=\mathbf{1}[q_j<0.05],\qquad
\rho=\frac{1}{d}\sum_{j=1}^{d}I_j,\qquad
A=\mathbf{1}[\rho\geq\tau_D].
\label{eq:drift}
\end{equation}
Here $\rho$ is a proportion of flagged features, not a probability of model failure. Marginal tests need not detect changes in feature dependence. BH is a defined comparison rule here, not a guarantee of dataset-level false-alert control under dependent, tied network-flow features.

\paragraph{Review boundary.}
We distinguish readiness, monitoring evidence (window, settings, scores, alert), and reviewed actions. An alert can request inspection while a model remains deployed. With labels, macro-F1, attack recall, and benign false-positive rate (FPR) provide separate evidence. Neither $A=1$ nor $A=0$ alone authorizes training, promotion, or rollback. This review boundary is a design recommendation, not an evaluated retraining policy or verified state machine.

\section{Controlled NIDS Experiment}
\label{sec:protocol}
\subsection{Provenance and Disjoint Pools}
CICIDS2017 provides labeled network flows~\cite{cicids}. Our input is a team-subsampled derivative of the third-party \emph{CICIDS2017: Cleaned \& Preprocessed} dataset, followed by the team's \emph{MLOps NIDS Dataset}.\footnote{\raggedright Kaggle source: \url{https://www.kaggle.com/datasets/ericanacletoribeiro/cicids2017-cleaned-and-preprocessed}; team derivative: \url{https://www.kaggle.com/datasets/hongtrnnguynvit/nids-reference-data}.} The local extraction script caps each label, reserves test rows, and shuffles the remaining rows. The experimental input, named \texttt{reference\_data.csv} by that earlier pipeline, contains 209,865 rows: 100,000 BENIGN, 50,000 DDoS, 50,000 PortScan, 8,150 BruteForce, and 1,715 WebAttacks, with 52 numeric features.

The associated notebook combines data, cleans values and features, groups labels, and selects features before our split. The local artifact lacks time or capture identifiers; we claim neither a temporal split nor removal of upstream feature-selection leakage. Complete byte-level derivation from the original release remains unverified.

We hash ordered numeric feature vectors, excluding labels and normalizing signed zero. Twenty rows in ten feature-identical, conflicting-label groups and one further duplicate are removed; there are no nonfinite rows. The remaining 209,844 rows are stratified into train/development/reference/evaluation pools at 50/10/20/20\% with split seed 260926. All six pool pairs have zero exact feature-vector overlap, including benign rows. This does not exclude session-level or correlated-flow dependence. The original input files remain unchanged.

\subsection{Fixed Model and Window Construction}
One binary XGBoost classifier~\cite{xgboost} is fitted on 74,994 BENIGN/DDoS rows from the train pool. PortScan, BruteForce, and WebAttacks are deliberately unseen during fitting. We use 200 trees, depth 6, learning rate 0.1, full row/column sampling, the histogram tree method, model seed 42, and prediction threshold 0.5. No scaler, SMOTE, or evaluation-based tuning is applied. Two development windows check execution only; the model remains fixed throughout evaluation.

A separate 5,000-row reference contains 70\% BENIGN and 30\% DDoS (reference sampling seed 260927). Evaluation windows contain $n\in\{500,1000\}$ rows sampled without replacement within each window. Table~\ref{tab:scenarios} specifies the scenarios. In S2, replacement attacks are divided as evenly as possible among the three held-out families. Thus S2 at 25\% replacement and $n=1000$ contains 700 benign, 225 DDoS, and 75 held-out-family rows. Replacement percentage is a fraction of \emph{attack rows}, not of the entire window.

\begin{table}[t]
\caption{Controlled scenarios; S2 changes attack composition at fixed prevalence.}
\label{tab:scenarios}
\centering
\begin{tabular}{@{}lll@{}}
\toprule
Scenario & Total attack fraction & Held-out share of attack rows \\
\midrule
S0 & 30\% & 0\% (all DDoS) \\
S1 & 15\%, 50\%, 85\% & 0\% (all DDoS) \\
S2 & 30\% & 25\%, 50\%, 100\% \\
\bottomrule
\end{tabular}
\end{table}

Each scenario/size uses seeds 42, 123, 456, 789, and 1024. S0 adds seeds 2000--2024, giving 30 stable windows per size and 120 windows overall. Windows may reuse rows from the same evaluation pool; standard deviations describe sampling with one model and one reference, not independent captures or training runs. Class counts are fixed by construction, including in S0.

\subsection{Detection, Metrics, and Reproducibility}
Two-sided asymptotic Kolmogorov--Smirnov tests cover all 52 features (6,240 tests), followed by Eq.~\ref{eq:drift}. We compare thresholds 0.10, 0.30, 0.50, and 0.70, retaining 0.30 as the principal setting and uncorrected tests as a comparison. Each window yields macro-F1, attack recall, benign FPR, ROC-AUC, family recall, and KS statistics. No threshold is selected retrospectively.

The artifact records settings, versions, checksums, exclusions, pool/window membership, tests, and metrics. An independent verifier reloads the model, recomputes predictions and all KS tests, and checks BH against SciPy; all checks passed. Versions include XGBoost 3.2.0, SciPy 1.17.1, and scikit-learn 1.8.0. Evidently's deployed preset does not explicitly configure this KS/BH rule, so the checks do not validate that adapter. Reproduction requires the checksummed local input; public versioned acquisition remains to be finalized.

\section{Results}
\label{sec:results}
\subsection{Distribution Drift and Predictive Quality}
Table~\ref{tab:results} reports the five principal windows per scenario at $n=1000$. S1 changes prevalence while retaining the learned attack family. Despite drift shares of 0.7423--0.9808, macro-F1 remains at least 0.9992. These alerts correctly reflect changed input composition; unchanged predictive quality does not make them false distribution alerts.

\begin{table}[t]
\caption{Fixed-model results for 1,000-row windows: mean$\pm$sample SD over five seeds. Alerts use $\tau_D=0.30$; S0 alert count uses all 30 stable windows.}
\label{tab:results}
\centering
\small
\begin{tabular*}{\linewidth}{@{\extracolsep{\fill}}lcccc@{}}
\toprule
Scenario & Drift share & Macro-F1 & Attack recall & Alerts \\
\midrule
S0 & 0.000$\pm$0.000 & 1.000$\pm$0.000 & 1.000$\pm$0.000 & 0/30 \\
S1 (15\%) & 0.742$\pm$0.059 & 1.000$\pm$0.000 & 1.000$\pm$0.000 & 5/5 \\
S1 (50\%) & 0.850$\pm$0.077 & 1.000$\pm$0.000 & 1.000$\pm$0.000 & 5/5 \\
S1 (85\%) & 0.981$\pm$0.000 & 0.999$\pm$0.001 & 1.000$\pm$0.001 & 5/5 \\
S2 (25\%) & 0.165$\pm$0.202 & 0.905$\pm$0.002 & 0.753$\pm$0.004 & 1/5 \\
S2 (50\%) & 0.865$\pm$0.073 & 0.789$\pm$0.003 & 0.508$\pm$0.005 & 5/5 \\
S2 (100\%) & 0.981$\pm$0.000 & 0.422$\pm$0.007 & 0.010$\pm$0.007 & 5/5 \\
\bottomrule
\end{tabular*}
\end{table}

S2 instead tests held-out attack families. At 25\% replacement, attack recall falls to 0.7533, while the principal detector alerts in only 1/5 windows. The same 1/5 count occurs at $n=500$, where mean recall is 0.7520. Hence low alert frequency can coexist with a material failure to recognize attacks in this constructed setting. At full replacement, recall is 0.0100 for $n=1000$: PortScan and WebAttacks recall are zero and BruteForce recall averages 0.0300. This is a limitation of the deliberately restricted BENIGN/DDoS baseline, not an evaluation of a model trained on all families.

The S2 trend must be interpreted as a mixture effect: for replacement fraction $r$, aggregate attack recall equals $(1-r)R_{\mathrm{DDoS}}+rR_{\mathrm{held\mbox{-}out}}$. Post-hoc inspection of the full evaluation pool finds 9,999/10,000 DDoS, 4/9,998 PortScan, 60/1,630 BruteForce, and 0/344 WebAttacks predicted as attacks. Increasing poorly recognized families therefore explains the decline; it does not demonstrate loss of DDoS knowledge over time. This inspection reuses the evaluation pool and is not an independent test.

S1 at 85\% prevalence and S2 at full replacement both flag approximately 98\% of features, yet their macro-F1 values differ greatly. Drift share thus cannot rank predictive harm in these cases. Across all 30 S0 windows at $n=1000$, mean macro-F1 is 0.999881, recall 0.999889, and FPR 0.000095; perfect values in the five-window table do not hold for every stable draw. S1 at 85\% prevalence has mean benign FPR 0.001333 in the principal windows; the other table rows have zero observed FPR.

\subsection{Threshold and Window-Size Sensitivity}
Table~\ref{tab:sensitivity} shows alert counts, with 30 windows for each S0 size and five for each shifted setting. All 60 S0 windows have zero BH alerts across the tested thresholds; this finite, conditional observation is not a zero-false-alarm guarantee. At 25\% replacement and $n=1000$, lowering $\tau_D$ from 0.30 to 0.10 increases alerts from 1/5 to 2/5, but does not eliminate missed shifts. At 50\% replacement and $\tau_D=0.70$, alerts increase from 1/5 for $n=500$ to 5/5 for $n=1000$.

\begin{table}[t]
\centering
\caption{BH-corrected alert counts across drift-share thresholds. S1 labels give total attack prevalence; S2 labels give the fraction of attack rows replaced by held-out families. S0 uses 30 windows per size; other scenarios use five.}
\label{tab:sensitivity}
\scriptsize
\begin{tabular*}{\linewidth}{@{\extracolsep{\fill}}l*{8}{c}@{}}
\toprule
& \multicolumn{4}{c}{$n=500$} & \multicolumn{4}{c}{$n=1000$} \\
\cmidrule(lr){2-5}\cmidrule(l){6-9}
Scenario & .10 & .30 & .50 & .70 & .10 & .30 & .50 & .70 \\
\midrule
S0 & 0/30 & 0/30 & 0/30 & 0/30 & 0/30 & 0/30 & 0/30 & 0/30 \\
S1 (15\%) & 5/5 & 5/5 & 5/5 & 3/5 & 5/5 & 5/5 & 5/5 & 3/5 \\
S1 (50\%) & 5/5 & 5/5 & 5/5 & 4/5 & 5/5 & 5/5 & 5/5 & 5/5 \\
S1 (85\%) & 5/5 & 5/5 & 5/5 & 5/5 & 5/5 & 5/5 & 5/5 & 5/5 \\
S2 (25\%) & 1/5 & 1/5 & 0/5 & 0/5 & 2/5 & 1/5 & 1/5 & 0/5 \\
S2 (50\%) & 5/5 & 5/5 & 4/5 & 1/5 & 5/5 & 5/5 & 5/5 & 5/5 \\
S2 (100\%) & 5/5 & 5/5 & 5/5 & 5/5 & 5/5 & 5/5 & 5/5 & 5/5 \\
\bottomrule
\end{tabular*}
\end{table}

Correction also changes sensitivity. Without BH, S0 at $n=1000$ and $\tau_D=0.10$ produces 3/30 alerts instead of 0/30. Conversely, S2 at 25\% replacement and $\tau_D=0.30$ produces 5/5 alerts instead of 1/5. This trade-off should be reported rather than choosing a test after observing its outcomes. Window size affects statistical power; these non-temporal windows do not measure detection delay.

\subsection{Serving Evidence Boundary}
The earlier bounded prototype replay used 100 Locust users for approximately 89 seconds, issuing 4,343 requests: 47.6 requests/s on average, 1.74\,s mean latency, 2.30\,s p95, and four HTTP 502 responses. It was not repeated for this revision. The available observation neither diagnoses those errors nor varies tenant/model counts; it cannot support scalability, availability, or comparative serving claims. The offline results above should not be interpreted as end-to-end validation of that platform.

\section{Discussion and Limitations}
\label{sec:discussion}
\paragraph{Implication for lifecycle design.}
Distribution alerts and labeled performance should remain separate, version-associated records. An absent alert should not suppress review when labels reveal failures. The registry makes this distinction explicit; our contribution is the scoped design and evaluation, not novelty in KS, BH, XGBoost, or registry technology.

\paragraph{Data and statistical validity.}
The study uses a preprocessed, subsampled CICIDS2017 derivative with upstream feature selection and no retained temporal metadata. Exact deduplication does not resolve all leakage or capture dependence. High scores on familiar classes may reflect dataset separability. Five shifted windows per setting, shared evaluation pools, fixed class counts, one model/reference, correlated features, and tied-value asymptotic tests limit inference. The construction measures changed composition and observed degradation; it does not isolate a causal mechanism or establish concept drift in $P(Y\mid X)$.

\paragraph{Security and operational scope.}
Tenant/model ownership checks address access to records, but do not sandbox model code. A malicious uploaded artifact can execute code through unsafe deserialization, dependency loading, or execution; containerization alone is not evidence of isolation. Restricted artifact formats, isolated execution identities, network/resource limits, and adversarial cross-tenant tests are required before accepting untrusted workloads. No such isolation evaluation is claimed here. Likewise, monitoring completeness under broker failures, model-cache behavior, multi-model/tenant load sweeps, alternative serving platforms, lifecycle ablations, and promotion/rollback gates remain unevaluated.

Further work needs temporal data, other environments, repeated model fits and reference samples, and retraining comparisons. The results support review, not a production threshold or continuous-training benefit.

\section{Conclusion}
A Native Registry links tenant/model serving and monitoring evidence. Controlled NIDS shifts show both strong drift without substantial degradation and declining recall with infrequent alerts. Sensitivity results reinforce separating alerts from update decisions. Natural drift, automated adaptation, and scalable secure serving require further evidence.

\begin{thebibliography}{99}

\bibitem{kreuzberger2023mlops}
D. Kreuzberger, N. K\"uhl, and S. Hirschl, ``Machine Learning Operations (MLOps): Overview, Definition, and Architecture,'' \textit{IEEE Access}, vol. 11, pp. 31866--31879, 2023. \url{https://doi.org/10.1109/ACCESS.2023.3262138}.

\bibitem{sculley2015debt}
D. Sculley, G. Holt, D. Golovin, E. Davydov, T. Phillips, D. Ebner,
V. Chaudhary, M. Young, J.-F. Crespo, and D. Dennison,
``Hidden Technical Debt in Machine Learning Systems,''
in \textit{Advances in Neural Information Processing Systems}, vol. 28, 2015.

\bibitem{amershi2019software}
S. Amershi, A. Begel, C. Bird, R. DeLine, H. Gall, E. Kamar,
N. Nagappan, B. Nushi, and T. Zimmermann,
``Software Engineering for Machine Learning: A Case Study,''
in \textit{ICSE-SEIP}, 2019, pp. 291--300.

\bibitem{subramanya2022devops}
R. Subramanya, S. Sierla, and V. Vyatkin, ``From DevOps to MLOps: Overview and Application to Electricity Market Forecasting,'' \textit{Applied Sciences}, vol. 12, no. 19, article 9851, 2022. \url{https://doi.org/10.3390/app12199851}.

\bibitem{renggli2021data}
C. Renggli, L. Rimanic, N. M. G\"urel, B. Karla\v{s}, W. Wu, and C. Zhang, ``A Data Quality-Driven View of MLOps,'' arXiv:2102.07750v1, 2021.

\bibitem{baylor2017tfx}
D. Baylor et al.,
``TFX: A TensorFlow-Based Production-Scale Machine Learning Platform,''
in \textit{KDD}, 2017, pp. 1387--1395.

\bibitem{mlflow}
M. Zaharia et al., ``Accelerating the Machine Learning Lifecycle with MLflow,'' \textit{IEEE Data Engineering Bulletin}, vol. 41, no. 4, pp. 39--45, 2018.

\bibitem{crankshaw2017clipper}
D. Crankshaw, X. Wang, G. Zhou, M. J. Franklin, J. E. Gonzalez, and I. Stoica,
``Clipper: A Low-Latency Online Prediction Serving System,''
in \textit{USENIX NSDI}, 2017, pp. 613--627.

\bibitem{Bhatt2025HarmonE}
H. Bhatt, S. Biswas, S. Rakhunathan, and K. Vaidhyanathan, ``HarmonE: A Self-Adaptive Approach to Architecting Sustainable MLOps,'' arXiv:2505.13693v1, 2025.

\bibitem{MarcosMercade2026MLOpsFrameworks}
J. Marcos-Mercad\'e, U. Lopez-Novoa, and M. Ega\~na Aranguren, ``An Empirical Evaluation of Modern MLOps Frameworks,'' arXiv:2601.20415v1, 2026.

\bibitem{gama2014survey}
J. Gama, I. \v{Z}liobait\.{e}, A. Bifet, M. Pechenizkiy, and A. Bouchachia, ``A Survey on Concept Drift Adaptation,'' \textit{ACM Computing Surveys}, vol. 46, no. 4, article 44, pp. 1--37, 2014. \url{https://doi.org/10.1145/2523813}.

\bibitem{webb2016drift}
G. I. Webb, R. Hyde, H. Cao, H.-L. Nguyen, and F. Petitjean,
``Characterizing Concept Drift,''
\textit{Data Mining and Knowledge Discovery}, vol. 30, no. 4,
pp. 964--994, 2016.

\bibitem{Zhang2025SSFNIDS}
X. Zhang, R. Zhao, Z. Jiang, H. Chen, Y. Ding, E. C. H. Ngai, and S.-H. Yang, ``Continual Learning with Strategic Selection and Forgetting for Network Intrusion Detection,'' arXiv:2412.16264v4, 2025.

\bibitem{xgboost}
T. Chen and C. Guestrin, ``XGBoost: A Scalable Tree Boosting System,'' in \textit{KDD}, 2016, pp. 785--794.

\bibitem{cicids}
I. Sharafaldin, A. H. Lashkari, and A. A. Ghorbani, ``Toward Generating a New Intrusion Detection Dataset and Intrusion Traffic Characterization,'' in \textit{ICISSP}, 2018, pp. 108--116.

\bibitem{sommer2010closedworld}
R. Sommer and V. Paxson,
``Outside the Closed World: On Using Machine Learning for Network Intrusion Detection,''
in \textit{Proc. IEEE Symposium on Security and Privacy}, 2010, pp. 305--316.

\bibitem{bh1995}
Y. Benjamini and Y. Hochberg, ``Controlling the False Discovery Rate: A Practical and Powerful Approach to Multiple Testing,'' \textit{Journal of the Royal Statistical Society: Series B}, vol. 57, no. 1, pp. 289--300, 1995. \url{https://doi.org/10.1111/j.2517-6161.1995.tb02031.x}.

\end{thebibliography}
\end{document}
